% Section I: Introduction -- Phase 1 in one paragraph, the gap, what the brief
% asks for in Phase 2, and our three contributions. Phase 1 numbers:
% results/metrics/test_results.csv (Random Forest 7.42 m, XGBoost 95.2% floor)
% and feature_check.csv (AP-affinity 15.87 -> 11.38 m, user-group CV).
\section{Introduction}\label{sec:intro}

A phone hears the Wi-Fi access points (WAPs) around it with different signal
strengths (RSSI, in dBm), and this list of readings, one \emph{scan}, works
like a fingerprint of the place where it was taken~\cite{he2016wifi}. On the
public UJIIndoorLoc dataset~\cite{torres2014ujiindoorloc} we predict three
things for every scan: its building and floor (classification) and its
position (regression, in meters).

In Phase~1 we solved the three tasks with classical machine learning: every
method used three separate models. The best, a Random Forest with our own
AP-affinity features, reached a mean 2D error of 7.42~m on the test set, while
XGBoost found the most floors (95.2\%). AP-affinity, which learns where each
WAP is usually heard, gave our largest gain (cross-validated error from 15.87
to 11.38~m). But the three task models share nothing, although building, floor
and position describe one place. Multi-task learning~\cite{caruana1997multitask}
trains related tasks together, so that what one task learns can help the
others.

Phase~2 builds the network the project brief asks for: one multilayer
perceptron (MLP) whose shared hidden layers feed three heads, two for
classification (building, floor) and one for regression (position). We keep the
protocol of Phase~1 unchanged: the same cleaned training data, every tuned
setting chosen by 5-fold user-group cross-validation (CV), which holds out whole
users and their phones, and the same test set, used only at the end. Our
contributions are:
\begin{itemize}
\item \emph{Two hybrids of Phase~1 and Phase~2.} Besides the scaled readings
  and heard (missingness) indicators of the WAPs, the network receives our
  Phase~1 features, the signal statistics and AP-affinity, as a ready-made
  location hint; and a blend averages the class probabilities and positions of
  the network and of the Phase~1 Random Forest, with one weight per task chosen
  by user-group CV.
\item \emph{An ablation study} that removes one part at a time (the shared
  trunk, the Phase~1 features, the heard indicators, the hidden WAPs, dropout
  and weight decay), with the CV noise of the random seed measured.
\item \emph{A fair comparison} of the network with our Phase~1 models and
  with single-task networks, on the same folds, test scans and metrics,
  including a hierarchical error breakdown, floor confusion matrices, error
  maps (in notebook~05) and a split into seen and new phones.
\end{itemize}
